\scanpage{6}
% archive foliation: 5
écrivant un élém de $E\otimes\widehat S$ \inserted{loc sur $S$} comme somme d'élts de la forme
\[
\omega\otimes f+e\otimes g\qquad
\begin{gathered}\omega\in\Omega(S),\ f,g\in\widehat S(S)\\e\text{ la section marquée de }\mathcal{O}_S\subset E\end{gathered}
\]
\[
\begin{aligned}
u_{a*}^{\widehat E}(\omega\otimes f+e\otimes g)
&=u_a(\omega)\otimes u_{a*}(f)+u_a(\unsure{\omega|e})\otimes u_{a*}(g)\\
&=\omega\otimes u_{a*}(f)+(e+a)\otimes u_{a*}(g)\\
&=(\omega\otimes u_{a*}(f)+a\otimes u_{a*}(g))+e\otimes u_{a*}(g)
\end{aligned}
\]
On a donc
\[
\begin{aligned}
\pi(u_{a*}^{\widehat E}(\omega\otimes f+e\otimes g))
&=\omega u_{a*}(f)+a u_{a*}(g)+u_{a*}(g)\\
&=\omega u_{a*}(f)+(a+1)u_{a*}(g)
\end{aligned}
\]
tandis que
\[
\pi(\omega\otimes f+e\otimes g)=\omega f+g
\]
Il faut avoir
\diagram{\begin{tikzcd}
E\otimes\widehat S\arrow[r,"\pi"]\arrow[d,"u_{a*}^{\widehat E}"']&\widehat S\arrow[d,"u_{a*}^{\widehat L}"]\\
E\otimes\widehat S\arrow[r,"\pi"]&\widehat S
\end{tikzcd}}
\unsure{Or}, $\pi\circ u_{a*}^{\widehat E}=u_{a*}^{\widehat L}\circ\pi$, et faisant $\unsure{\omega\otimes f=0|\omega f=0}$, on trouve eq.
\[
\boxed{u_{a*}^{\widehat L}(g)=(a+1)u_{a*}(g)}
\quad\left(=\frac{g}{(1+a)^{n-1}}\text{ si }\deg g=n\right)
\]
Vérifions \unsure{ceci} en général par acquit de conscience.
\[
\begin{aligned}
u_{a*}^{\widehat L}\pi(\omega\otimes f+e\otimes g)
&=(a+1)u_{a*}(\omega f+g)\\
&=(a+1)(\underbrace{u_{a*}(\omega)}_{\frac{\omega}{\unsure{1-a|1+a}}}u_{a*}(f)+u_{a*}(g))\\
&=\\
&=\omega u_{a*}(f)+(a+1)u_{a*}(g)\qquad\text{OK !}
\end{aligned}
\]
Si on revient au cas général d'une extension quelc. de $\mathcal{O}_S$ par $\Omega$, définie par un torseur $\mathcal{T}$ sous $\Omega$, on trouve donc
\[
\left\{\begin{aligned}
\widehat S&=\deleted{\unsure{\widehat{\operatorname{Sym}}(\Omega)\wedge_{(\Omega,u)}}}\mathcal{T}\wedge_{(\Omega,u_{a*})}\widehat{\operatorname{Sym}}(\Omega)\\
\widehat L&=\mathcal{T}\wedge_{(\Omega,u_{a*}^{\widehat L})}\widehat{\operatorname{Sym}}(\Omega).
\end{aligned}\right.
\]
\deleted{Comment calculer les extensions successives $\mathcal{J}^n/\mathcal{J}^{n+2}$ ext. de $\operatorname{Sym}^n\Omega$ par $\operatorname{Sym}^{n+1}\Omega$}
